\structelem{ПЕРЕЧЕНЬ СОКРАЩЕНИЙ И ОБОЗНАЧЕНИЙ}

В настоящем отчете о НИР применяют следующие сокращения и
обозначения.

\vspace{6pt}

% Перечень сокращений по ГОСТ 7.32-2017, п. 6.15:
% слева без абзацного отступа --- сокращение, справа через тире ---
% его расшифровка, без знаков препинания в конце строки.
\begin{description}[
    font=\normalfont\bfseries,
    leftmargin=0pt, labelindent=0pt, labelsep=0.4em,
    style=sameline, itemsep=2pt, topsep=0pt
]
    \item[AE] {\normalfont --- Autoencoder, автоэнкодер}
    \item[ATT\&CK] {\normalfont --- Adversarial Tactics, Techniques and Common Knowledge --- матрица MITRE}
    \item[BiLSTM] {\normalfont --- Bidirectional Long Short-Term Memory}
    \item[cGAN] {\normalfont --- Conditional Generative Adversarial Network}
    \item[CI] {\normalfont --- Confidence Interval, доверительный интервал}
    \item[CNN] {\normalfont --- Convolutional Neural Network, сверточная нейронная сеть}
    \item[DL] {\normalfont --- Deep Learning, глубокое обучение}
    \item[DPI] {\normalfont --- Deep Packet Inspection, глубокий анализ пакета}
    \item[ECE] {\normalfont --- Expected Calibration Error}
    \item[EDA] {\normalfont --- Exploratory Data Analysis}
    \item[EPS] {\normalfont --- Events Per Second, событий в секунду}
    \item[FAR] {\normalfont --- False Alarm Rate}
    \item[FN] {\normalfont --- False Negative, ложноотрицательный}
    \item[FP] {\normalfont --- False Positive, ложноположительный}
    \item[FPR] {\normalfont --- False Positive Rate}
    \item[FSM] {\normalfont --- Finite-State Machine, конечный автомат}
    \item[GAN] {\normalfont --- Generative Adversarial Network}
    \item[GRU] {\normalfont --- Gated Recurrent Unit}
    \item[IAT] {\normalfont --- Inter-Arrival Time, межпакетный интервал}
    \item[IDS] {\normalfont --- Intrusion Detection System}
    \item[IPFIX] {\normalfont --- IP Flow Information Export}
    \item[IPS] {\normalfont --- Intrusion Prevention System}
    \item[IQR] {\normalfont --- Interquartile Range, межквартильный размах}
    \item[KDD] {\normalfont --- Knowledge Discovery in Databases}
    \item[KL] {\normalfont --- Kullback--Leibler divergence}
    \item[KS-test] {\normalfont --- Kolmogorov--Smirnov test}
    \item[LSTM] {\normalfont --- Long Short-Term Memory}
    \item[MCC] {\normalfont --- Matthews Correlation Coefficient}
    \item[ML] {\normalfont --- Machine Learning, машинное обучение}
    \item[MLP] {\normalfont --- Multilayer Perceptron, многослойный перцептрон}
    \item[MMD] {\normalfont --- Maximum Mean Discrepancy}
    \item[NIDS] {\normalfont --- Network-based Intrusion Detection System}
    \item[NTA] {\normalfont --- Network Traffic Analysis}
    \item[OCSVM] {\normalfont --- One-Class Support Vector Machine}
    \item[OOD] {\normalfont --- Out-Of-Distribution}
    \item[PR-AUC] {\normalfont --- Area Under the Precision--Recall Curve}
    \item[PSI] {\normalfont --- Population Stability Index}
    \item[ReLU] {\normalfont --- Rectified Linear Unit}
    \item[RF] {\normalfont --- Random Forest, случайный лес}
    \item[RNN] {\normalfont --- Recurrent Neural Network, рекуррентная нейронная сеть}
    \item[ROC-AUC] {\normalfont --- Area Under the Receiver Operating Characteristic Curve}
    \item[SOAR] {\normalfont --- Security Orchestration, Automation and Response}
    \item[SOC] {\normalfont --- Security Operations Center, центр мониторинга безопасности}
    \item[SVM] {\normalfont --- Support Vector Machine}
    \item[TCN] {\normalfont --- Temporal Convolutional Network}
    \item[TN] {\normalfont --- True Negative, истинноотрицательный}
    \item[TP] {\normalfont --- True Positive, истинноположительный}
    \item[TPR] {\normalfont --- True Positive Rate, чувствительность (Recall)}
    \item[TTD] {\normalfont --- Time-To-Detect, время до обнаружения}
    \item[VAE] {\normalfont --- Variational Autoencoder, вариационный автоэнкодер}
    \item[XGBoost] {\normalfont --- eXtreme Gradient Boosting}
\end{description}
